\section{The Plan}
\label{sec:ch07-the-plan}
\index{query plan|(}%

Between the document arriving and the first subgraph being called, the router
decides what it is going to do. That decision is the query plan, and Cosmo will
show it to you.

\index{query plan!X-WG-Include-Query-Plan header@\code{X-WG-Include-Query-Plan} header}%
Three request headers are involved~\autocite{cosmo2026queryplan}:

\begin{minted}{text}
X-WG-Include-Query-Plan: true
X-WG-Skip-Loader: true
X-WG-Disable-Tracing: true
\end{minted}

\medskip\noindent\textbf{Show me the plan.} The first puts it under
\code{extensions.queryPlan}, beside the data.

\medskip\noindent\textbf{Do not run it.} The second stops the router executing
the plan it just made. \code{data} comes back null and the subgraphs are never
called, which makes asking for a plan free.

\medskip\noindent\textbf{Keep it out of the traces.} The third keeps a plan
request out of the router's own tracing, so that reading a plan does not show up
in your metrics as traffic.

\medskip
The first two need the development mode of
section~\ref{sec:ch07-three-processes}; I did not test the third without it.
Send the chapter's query with both set and this comes back, whitespace collapsed
inside the embedded query strings and nothing else touched.

\begin{minted}[fontsize=\footnotesize]{json}
{
  "version": "1",
  "kind": "Sequence",
  "children": [
    {
      "kind": "Single",
      "fetch": {
        "kind": "Single",
        "subgraphName": "catalog",
        "subgraphId": "0",
        "fetchId": 0,
        "query": "{ products { title price __typename id } }"
      }
    },
    {
      "kind": "Single",
      "fetch": {
        "kind": "BatchEntity",
        "path": "products",
        "subgraphName": "reviews",
        "subgraphId": "1",
        "fetchId": 1,
        "dependsOnFetchIds": [0],
        "representations": [
          {
            "kind": "@key",
            "typeName": "Product",
            "fragment": "fragment Key on Product { __typename id }"
          }
        ],
        "query": "query($representations: [_Any!]!) { _entities(representations: $representations) { ... on Product { __typename reviews { rating body } } } }",
        "dependencies": [
          {
            "coordinate": { "typeName": "Product", "fieldName": "reviews" },
            "isUserRequested": true,
            "dependsOn": [
              {
                "fetchId": 0,
                "subgraph": "catalog",
                "coordinate": { "typeName": "Product", "fieldName": "id" },
                "isKey": true,
                "isRequires": false
              }
            ]
          }
        ]
      }
    }
  ],
  "normalizedQuery": "{products {title price reviews {rating body}}}"
}
\end{minted}

\index{query plan!Sequence node@\code{Sequence} node}%
The outer \code{kind} is \code{Sequence}, which is the plan saying these run in
order rather than at once. There is a \code{Parallel} node as well, and this
query does not get one, for a reason the plan itself gives a few lines further
down. Figure~\ref{fig:ch07-query-plan} is the same structure without the
punctuation.

\begin{figure}[htbp]
  \centering
  % The router's plan for { products { title price reviews { rating body } } },
% drawn from the JSON the router returns under X-WG-Include-Query-Plan.
% Referenced from chapter 07. Bare tikzpicture; the figure environment, caption
% and label live at the call site.
%
% The dependency arrow runs underneath the two fetch nodes rather than between
% them: its label is wider than the gap and overprinted the boxes.
\begin{tikzpicture}[
    font=\small,
    fetch/.style={draw, rounded corners=2pt, minimum width=38mm,
                  minimum height=13mm, align=center},
    root/.style={draw, rounded corners=2pt, minimum width=28mm,
                 minimum height=8mm, align=center, fill=black!8},
    tree/.style={-{Stealth[length=2mm]}},
    dep/.style={-{Stealth[length=2mm]}, dashed},
    lbl/.style={font=\scriptsize, align=center, fill=white, inner sep=2pt},
    note/.style={font=\scriptsize\itshape, align=center}
  ]

  \node[root] (seq) at (0,0) {\code{Sequence}};

  \node[fetch] (f0) at (-3.6,-2.0)
    {\code{fetch 0}\\ \code{Single}\\ subgraph \code{catalog}};
  \node[fetch] (f1) at (3.6,-2.0)
    {\code{fetch 1}\\ \code{BatchEntity}\\ subgraph \code{reviews}};

  \draw[tree] (seq.south) -- (-3.6,-0.75) -- (f0.north);
  \draw[tree] (seq.south) -- (3.6,-0.75) -- (f1.north);

  % The dependency is what forces the order. Everything else is structure.
  \draw[dep] (-3.6,-2.65) -- (-3.6,-3.5) -- (3.6,-3.5) -- (3.6,-2.65);
  \node[lbl] at (0,-3.5)
    {\code{dependsOnFetchIds: [0]}\\
     \code{Product.id}, \code{isKey: true}};

  \node[note, anchor=north, text width=115mm] at (0,-4.1)
    {Solid lines are the plan's own structure. The dashed line is the data
     dependency underneath it: fetch 1 cannot be sent until fetch 0 has
     answered, because the key it needs is in that answer.};

\end{tikzpicture}

  \caption{Two fetches under one \code{Sequence}. The tree says what runs; the
  \code{dependencies} array says why it cannot run any other way. Finding the
  dashed line is how you count a query's round trips before running it, which
  is not the same as knowing what they cost.}
  \label{fig:ch07-query-plan}
\end{figure}

\index{query plan!BatchEntity fetch@\code{BatchEntity} fetch}%
Fetch 0 is a \code{Single}: an ordinary query, sent to one subgraph, depending
on nothing. Fetch 1 is a \code{BatchEntity}, which is the plan's name for an
\code{\_entities} call, and the \code{path} field says where in the response its
results belong. \code{dependsOnFetchIds: [0]} is the ordering constraint, and
\code{representations} describes what fetch 1 will send: for every
\code{Product} at that path, \code{\_\_typename} and \code{id}, written as a
fragment.

\index{query plan!dependencies array@\code{dependencies} array}%
The \code{dependencies} array is the part I did not expect and now look at
first. It is the planner showing its work: the client asked for
\code{Product.reviews}, so \code{isUserRequested} is true, and answering it
depends on \code{Product.id} from subgraph \code{catalog}, which is a key rather
than a \code{@requires}. Two booleans and a coordinate, and between them they
say exactly why there are two round trips instead of one. When a query is
slower than I expected, this array is where I start.

\subsection*{Where \code{\_\_typename} and \code{id} came from}

Look at what fetch 0 sends:

\begin{minted}{graphql}
{ products { title price __typename id } }
\end{minted}

The client asked for \code{title} and \code{price}. It did not ask for
\code{\_\_typename} and it did not ask for \code{id}. The router added both,
because fetch 1 cannot be built without them, and the fragment under
\code{representations} is the shopping list.

That is easy to state and easy to half-believe, so here is the control. Same
root field, same subgraph, one selection removed:

\begin{minted}{graphql}
{ products { title price } }
\end{minted}

\begin{minted}[fontsize=\footnotesize]{json}
{
  "kind": "Sequence",
  "children": [
    { "kind": "Single",
      "fetch": { "kind": "Single", "subgraphName": "catalog", "fetchId": 0,
                 "query": "{ products { title price } }" } }
  ],
  "normalizedQuery": "{products {title price}}"
}
\end{minted}

One child, and no key. Nothing crosses a boundary, so the router fetches nothing
to cross it with. The extra fields are not a habit the router has; they are a
dependency it worked out.

The mirror image makes the same point from the other side. Ask for reviews and
nothing else:

\begin{minted}{graphql}
{ products { reviews { rating } } }
\end{minted}

Catalog is still called, and called for nothing whatsoever that the client
wanted:

\begin{minted}[fontsize=\footnotesize]{text}
Single       catalog   { products { __typename id } }
BatchEntity  reviews   query($representations: [_Any!]!) { _entities(...) }
\end{minted}

A whole subgraph round trip to learn three identifiers. There is no way around
it: \code{products} is Catalog's field, and the list of which products exist is
Catalog's answer to give. This is the shape that makes people ask whether
federation is worth it, and
chapter~\ref{ch:do-you-actually-need-federation} takes the question seriously.
What is worth saying here is that the cost is visible before you deploy
anything. The plan is a build-time artefact of a query, and you can read it in
Postman.

\subsection*{When the router does not wait}

\index{query plan!Parallel node@\code{Parallel} node}%
The \code{Sequence} above is not the router being cautious. Ask for two root
fields that owe each other nothing:

\begin{minted}{graphql}
{ products { title } reviews { rating } }
\end{minted}

and the two fetches land under a different node:

\begin{minted}[fontsize=\footnotesize]{json}
{
  "kind": "Sequence",
  "children": [
    { "kind": "Parallel",
      "children": [
        { "kind": "Single",
          "fetch": { "subgraphName": "catalog", "fetchId": 0,
                     "query": "{ products { title } }" } },
        { "kind": "Single",
          "fetch": { "subgraphName": "reviews", "fetchId": 1,
                     "query": "{ reviews { rating } }" } }
      ] }
  ]
}
\end{minted}

Neither fetch has a \code{dependsOnFetchIds}, so neither has to wait, and both
go at once. That is the rule the two plans share: the router serialises exactly
as much as the data forces it to, and a \code{Sequence} of two children means it
found a reason.

\subsection*{The plan is not the traffic}

One caution before the next section. Everything above is the router's account of
what it intends to do, printed by the router. It is not evidence that any of it
happened, and with \code{X-WG-Skip-Loader} set it is evidence that none of it
did.

Section~\ref{sec:ch07-the-traffic-in-full} reads the other end: what the
subgraphs actually received. The two agree in this sample, and I only know that
because the repository's verification script checks both and would fail if they
stopped agreeing.

\index{query plan|)}%
